\documentclass[11pt]{article}
\usepackage[margin=1in]{geometry}
\usepackage{amsmath,amssymb,amsthm}
\usepackage{microtype}
\usepackage{enumitem}
\usepackage[colorlinks=true,linkcolor=blue,citecolor=blue,urlcolor=blue]{hyperref}
\hypersetup{pdftitle={A constructive characterization of exactly trackable heat fluxes},
 pdfauthor={Kaiwen Jin},pdfsubject={Exact range of separated boundary heat flux tracking}}
\newtheorem{theorem}{Theorem}[section]
\newtheorem{lemma}[theorem]{Lemma}
\newtheorem{proposition}[theorem]{Proposition}
\theoremstyle{remark}
\newtheorem{remark}[theorem]{Remark}
\newcommand{\R}{\mathbb R}
\newcommand{\N}{\mathbb N}
\newcommand{\Om}{\Omega}
\newcommand{\dt}{\partial_t}
\newcommand{\dn}{\partial_\nu}
\newcommand{\calX}{\mathcal X}
\newcommand{\calR}{\mathcal R}
\newcommand{\calF}{\mathcal F}
\newcommand{\calY}{\mathcal Y}
\DeclareMathOperator{\dist}{dist}
\title{A constructive characterization of exactly trackable heat fluxes}
\author{Kaiwen Jin\\
\small\href{mailto:yui.ui.siki@gmail.com}{yui.ui.siki@gmail.com}}
\date{12 October 2026}
\begin{document}
\maketitle
\begin{abstract}
We characterize the normal heat fluxes that can be produced on an observation
patch by square-integrable Dirichlet data on a disjoint control patch.
The domain is any bounded connected smooth Euclidean domain in dimension
at least two. The characterization consists of explicit conditions on the
target: causal smoothness, weighted Taylor coefficients computed from a
Gaussian surface potential, agreement of the resulting series, and
volume, boundary-trace, support and observation compatibility tests.
No control or caloric extension is an unknown in these tests.
The reconstructed parallel boundary values define the exact causal
control norm. In this norm the squared integral of the separated
heat-flux kernel gives the exact norm of point evaluation at each positive
time. Every fixed target is little-o of this scale at the initial time,
and no prescribed further decay factor holds for all targets.
\end{abstract}

\section{Introduction}

A boundary heater prescribes temperature on one part of a body, whereas a
sensor measures heat flux on another. Exact tracking asks which complete
time-dependent flux profiles can be produced by such a heater. Diffusion
smooths a signal before it reaches a separated observation patch, but
smoothness alone does not describe the attainable profiles: continuation
through the body and support of the boundary input impose additional
compatibility conditions.

We consider a bounded connected $C^\infty$ domain $\Om\subset\R^d$,
$d\geq2$, a time $T>0$, and nonempty relatively open sets
$\Gamma,\Sigma\subset\partial\Om$ with
$\overline\Gamma\cap\overline\Sigma=\varnothing$.
The input is any $u\in L^2((0,T)\times\Gamma)$.
Let $y_u$ be the zero-initial transposition solution of
\begin{equation}\label{eq:state}
 \dt y_u-\Delta y_u=0,\qquad y_u(0,\cdot)=0,\qquad
 y_u|_{\partial\Om}=\mathbf1_\Gamma u.
\end{equation}
With the outward normal $\nu$, the tracking range is
\begin{equation}\label{eq:range}
 \calR_T=\{\dn y_u|_{(0,T)\times\Sigma}:
                  u\in L^2((0,T)\times\Gamma)\}.
\end{equation}
The separation of the closures makes each output smooth locally on
$[0,T]\times\Sigma$, including all one-sided derivatives at $T$,
and flat at $t=0$.
Throughout, a flux is identified with this smooth representative.
No regularity of the edges of the two patches is assumed.

B\'arcena-Petisco and Zuazua study tracking controllability for the heat
equation \cite{BZ}. Their one-dimensional sufficient construction admits
flat Gevrey targets of order below two; they also discuss approximation
and the problem of finding a sharp tracking class.
For interval geometry, Davron and Lissy give exact Gevrey-based tracking
criteria, with a reachable-state compatibility condition at the finite-time
endpoint. Their variants include Dirichlet input and Neumann output
\cite[Theorems~1--2 and Section~3.1.2]{DL}.
The related lateral Cauchy problem prescribes both temperature and normal
flux on part of the boundary.
Puzyrev and Shlapunov express its solvability through extension of a heat
potential to a caloric function \cite[Theorem~2]{PS}.
That criterion identifies the relevant continuation problem, but the
extension remains an unknown.
Here we retain the heat-potential construction and give explicit tests
for its continuation, together with the conditions needed for the
original $L^2$ input class.

More precisely, for every geometry and every pair of patches just described,
we define a class $\calX_T$ using only the target, known Gaussian integrals,
fixed binomial translations and limits of explicitly computed functions.
Theorem~\ref{thm:range} proves $\calR_T=\calX_T$.
The tests use one small seed patch, yet enforce compatibility on all of
$\Sigma$ and support on the exact original $\Gamma$.
They give a unique recovered input and a concrete norm $N_f(T)$, with
$N_f(t)$ equal to its $L^2$ size up to time $t$.
Theorem~\ref{thm:sharp} then gives the exact geometry-dependent scale of
point evaluation and its optimal initial-time restrictions.
This is an infinite regularity and compatibility criterion; no
finite decision procedure or universal Sobolev or Gevrey description is
asserted.

The main difficulty is to combine spatial continuation with rough boundary
data. We cancel the measured flux source by subtracting the state extended
by zero across the seed patch from a heat single-layer potential.
The difference is causal and caloric on an enlarged domain.
A separated-source Gaussian estimate supplies Taylor radii fixed by
geometry, uniformly up to both time endpoints.
The target therefore determines every continuation coefficient by
translation. Absolute coefficient sums and overlap equalities test whether
these local series form one caloric function.
Parallel boundary values then distinguish precisely the original
$L^2$ states, including inputs with no positive Sobolev regularity.
Their convergence is proved directly from heat-Poisson mass and
concentration estimates.

Section~\ref{sec:class} defines the target class and states the range theorem.
Section~\ref{sec:trace} proves the boundary trace facts,
Section~\ref{sec:continuation} proves the continuation estimates,
and Section~\ref{sec:rangeproof} establishes the equivalence and norm.
Section~\ref{sec:sharp} proves the sharp quantitative assertions.
We write $\N=\{0,1,\ldots\}$ and use the usual multi-index notation:
$|\alpha|=\sum_i\alpha_i$, $\alpha!=\prod_i\alpha_i!$,
$h^\alpha=\prod_i h_i^{\alpha_i}$, and $\alpha\geq\beta$ componentwise.
All statements hold for real data; conjugation in the proofs also gives the
complex version.

\section{Explicit conditions on the target}\label{sec:class}

We first fix auxiliary objects using the geometry alone.
Choose a ball $B$ centered at a point of $\Sigma$, small enough that
$\overline B\cap\partial\Om$ is a compact subset of $\Sigma$ and
\[
 E=B\setminus\overline\Om
\]
is connected. Choose $\chi\in C^\infty_c(\Sigma)$ equal to one on
$\overline B\cap\partial\Om$, and set $D=\Om\cup B$.
Thus $D$ is open, bounded and connected; its boundary need not be smooth.
The ball permits cancellation of the boundary source, and $\chi$ makes
the known potential use only a compact part of the observed flux.

List all rational points of $D$ as $\{q_j:j\in\N\}$, with a root
$q_0\in E$, and put
\begin{equation}\label{eq:radii}
 r_j=\frac{\dist(q_j,\partial D)}{128\sqrt d}.
\end{equation}
Fix a rooted tree on these centers with finite root-to-vertex paths and
\begin{equation}\label{eq:edges}
 |q_k-q_j|<r_j/4
 \quad\text{whenever $q_j$ is the parent of $q_k$.}
\end{equation}
These balls cover $D$. Finally fix a tubular radius $\varepsilon_0>0$
for $\partial\Om$ and set
$\varepsilon_n=\varepsilon_0\,2^{-n-2}$, $n\in\N$.
The existence of the geometric choices is justified below.

Here are the five tests defining $\calX_T$.
\begin{enumerate}[label=\textnormal{(\Alph*)},leftmargin=*]
\item\label{test:A}
The function $f$ belongs to $C^\infty_{\rm loc}([0,T]\times\Sigma)$
and every mixed space-time derivative vanishes at $t=0$.
Smoothness at $T$ means continuity of each derivative from the left on
compact spatial subsets.
\item\label{test:B}
Form the known potential
\begin{equation}\label{eq:potential}
 V_f(x,t)=\int_0^t\int_\Sigma
       \Phi(t-s,x-z)\chi(z)f(s,z)\,dS(z)\,ds,\qquad
 \Phi(r,w)=(4\pi r)^{-d/2}e^{-|w|^2/(4r)}.
\end{equation}
At the root set
$c^{(0)}_\alpha(t)=\partial_x^\alpha V_f(q_0,t)/\alpha!$.
For each tree edge in \eqref{eq:edges} compute
\begin{equation}\label{eq:translation}
 c^{(k)}_\beta(t)=
 \sum_{\alpha\geq\beta}\binom{\alpha}{\beta}
 c^{(j)}_\alpha(t)(q_k-q_j)^{\alpha-\beta}.
\end{equation}
Require these sums, with every time derivative, to converge uniformly on
$[0,T]$. The coefficients are computed recursively along the fixed finite
paths; none is an unknown chosen to satisfy a condition.
For every $j,m\in\N$ require
\begin{equation}\label{eq:budget}
 \sum_{\alpha\in\N^d}r_j^{|\alpha|}
      \sup_{0\leq t\leq T}|\partial_t^m c^{(j)}_\alpha(t)|<\infty.
\end{equation}
Also require, for every $j,\alpha,m$,
\begin{equation}\label{eq:recurrence}
 \partial_t^m c^{(j)}_\alpha(0)=0,\qquad
 \dt c^{(j)}_\alpha=
 \sum_{i=1}^d(\alpha_i+2)(\alpha_i+1)c^{(j)}_{\alpha+2e_i}.
\end{equation}
On each overlap of $B(q_j,r_j)$ and $B(q_k,r_k)$ require equality of
the two series at every rational space-time point in that overlap.
The summability conditions imply continuity, so this countable test gives equality
everywhere. If any sum diverges, the target fails this test.

When this test passes, define the explicitly computed functions
\begin{equation}\label{eq:FY}
 \calF_f(x,t)=\sum_\alpha c^{(j)}_\alpha(t)(x-q_j)^\alpha
       \quad(x\in B(q_j,r_j)),\qquad
 \calY_f=V_f-\calF_f\quad\text{on $\Om$.}
\end{equation}
The first definition is independent of the covering ball.
\item\label{test:C}
Require $\calY_f\in L^2((0,T)\times\Om)$.
Define its parallel boundary values, with the normal fixed at $p$,
\begin{equation}\label{eq:Gn}
 G_n(f)(t,p)=\calY_f(t,p-\varepsilon_n\nu(p)).
\end{equation}
Require the scalar Cauchy condition
\begin{equation}\label{eq:cauchy}
 \lim_{N\to\infty}\sup_{n,m\geq N}
 \int_0^T\int_{\partial\Om}|G_n(f)-G_m(f)|^2\,dS\,dt=0.
\end{equation}
\item\label{test:D}
Require exact control-patch compatibility:
\begin{equation}\label{eq:support}
 \lim_{n\to\infty}
 \int_0^T\int_{\partial\Om\setminus\Gamma}|G_n(f)|^2\,dS\,dt=0.
\end{equation}
The complement is literal. In particular this condition includes patch
edges of positive surface measure.
\item\label{test:E}
For every compact coordinate set $K\Subset\Sigma$ require
\begin{equation}\label{eq:observation}
 \nu(p)\cdot\nabla_x\calY_f(t,p-\varepsilon_n\nu(p))
          \longrightarrow f(t,p)
       \quad\text{in }C^\infty([0,T]\times K).
\end{equation}
This means uniform convergence of every derivative in time and local
boundary coordinates. The test uses the entire observation patch,
including components away from the seed ball.
\end{enumerate}
Let $\calX_T=\calX_T(\Om,\Gamma,\Sigma)$ be the space passing these tests.
All functions in the tests are given by \eqref{eq:potential}--\eqref{eq:FY};
the boundary limits are tests of those computed functions.

\begin{theorem}[Exact range]\label{thm:range}
For every $\Om,d,T,\Gamma,\Sigma$ satisfying the hypotheses of
Section~1, and every admissible choice of the auxiliary geometry above,
\[
 \calR_T=\calX_T.
\]
Each target in this space is generated by a unique
$u\in L^2((0,T)\times\Gamma)$. It is recovered by
$u=\lim_n G_n(f)|_\Gamma$ in $L^2$.
For $0\leq t\leq T$ the formula
\begin{equation}\label{eq:N}
 N_f(t)^2=\lim_{n\to\infty}
              \int_0^t\int_\Gamma|G_n(f)(s,p)|^2\,dS(p)\,ds
\end{equation}
equals $\|u\|_{L^2((0,t)\times\Gamma)}^2$.
The norm $\|f\|_{\calX_T}=N_f(T)$ makes the range a Hilbert space
isometric to $L^2((0,T)\times\Gamma)$.
The class and this norm do not depend on the auxiliary choices.
\end{theorem}

The theorem expresses membership without a singular system, a dual
inequality, or an existential control or extension.
The boundary limit yields the control after membership has been checked.
Condition \eqref{eq:observation} is needed because the initial continuation
uses only a small part of $\Sigma$; condition \eqref{eq:support} retains
the actual control patch.

\begin{lemma}[Geometric choices]\label{lem:geometry}
The seed ball, covering balls, tree and tubular coordinates used above exist.
For sufficiently small $\varepsilon>0$,
$\Psi_\varepsilon(p)=p-\varepsilon\nu(p)$ parametrizes
$\partial\Om_\varepsilon$, where
$\Om_\varepsilon=\{x\in\Om:\dist(x,\partial\Om)>\varepsilon\}$.
Its outward normal is $\nu(p)$ and its surface Jacobian
$J_\varepsilon(p)$ tends smoothly and uniformly to one.
There is a uniform interior and exterior tangent-ball radius.
\end{lemma}
\begin{proof}
Finite smooth graph charts on the compact boundary give, uniformly in
$p,z\in\partial\Om$,
\begin{equation}\label{eq:quadratic}
 |\nu(p)-\nu(z)|\leq C|p-z|,\qquad
 |(p-z)\cdot\nu(p)|+|(p-z)\cdot\nu(z)|\leq C|p-z|^2.
\end{equation}
For distant pairs enlarge $C$ using compactness.
The normal-coordinate derivative at distance zero is invertible.
The local inverse follows, for example, by shrinking a chart until the
equation has the form $v=b-H(v)$ with $\|DH\|<1/2$ and iterating;
differentiation gives a smooth inverse.
If $p-r\nu(p)=z-s\nu(z)$, taking the scalar product with $p-z$ in
\eqref{eq:quadratic} gives
$|p-z|^2\leq C(|r|+|s|)|p-z|^2$.
Thus small normal coordinates are globally injective.
A nearest boundary point exists; its displacement is normal and the graph
fixes the inward sign. The normal distance is therefore the true distance.
Tangential derivatives of $\nu$ are tangent, giving the asserted normal
and $J_\varepsilon=\det(I-\varepsilon\,d\nu)$.
The inner domain may have several components.
For small $R>0$,
\[
 |p-z-R\nu(z)|^2-R^2\geq(1-2CR)|p-z|^2.
\]
The ball centered at $z+R\nu(z)$ has no boundary point inside, and its
center is exterior by the normal chart. Connectedness of the ball makes
it exterior. Reversing the sign gives an interior tangent ball.

For the seed ball rotate a graph to $x_d=h(x')$, with
$h(0)=Dh(0)=0$. On a radius-$a$ ball let
$H(x')=(a^2-|x'|^2)^{1/2}$. Since $D^2H\leq-I/a$ and $D^2h$ is bounded,
$H-h$ is strictly concave for small $a$.
Its positive set is convex and each exterior vertical fiber is
$(\max(h,-H),H)$. Their continuous midpoint section proves that $E$
is connected. The same ball meets $\Om$, so $D$ is connected.

An open connected set is polygonally connected: the set of points reachable
from one point by interior polygonal paths is open and has open complement.
A fixed path has positive clearance from $\partial D$.
Dense rational points along a sufficiently fine subdivision therefore
give a directed chain satisfying \eqref{eq:edges}.
Every center is reachable. Choose a predecessor at the preceding minimal
integer path length, resolving ties by the fixed listing; this gives the tree.
If $\dist(x,\partial D)=a$, choose $q_j$ within $a/(512\sqrt d)$ of $x$.
Then $r_j>a/(256\sqrt d)>|x-q_j|$, proving the covering assertion.
\end{proof}

\section{The original \texorpdfstring{$L^2$}{L2} boundary class}\label{sec:trace}

For $g\in L^2((0,T)\times\partial\Om)$, the zero-initial transposition
identity is
\begin{equation}\label{eq:transposition}
 \int_0^T\int_\Om y(-\dt\phi-\Delta\phi)\,dx\,dt
       =-\int_0^T\int_{\partial\Om}g\,\dn\phi\,dS\,dt
\end{equation}
for every $\phi\in C^\infty([0,T]\times\overline\Om)$ with
$\phi|_{\partial\Om}=0$ and $\phi(T)=0$.
The state belongs to $L^2((0,T)\times\Om)$.
We prove the needed representation and trace facts with this identity.
We use monotone and dominated convergence, integrable Fubini,
H\"older's inequality and $L^2$ completeness in their usual forms
\cite[Theorems~4.6, 4.24, 5.18, 7.7, 7.10]{HM}.
Smooth-chart mollification and the spatial Fourier transform on
distributions use the conventions of
\cite[Sections~1.9--1.12, 5.A--5.B]{HP}.

\subsection{Heat-kernel mass and concentration}
Let $p(r,x,y)$ be the Dirichlet heat kernel of $\Om$, and set
\begin{equation}\label{eq:P}
 P(r,x,z)=-\partial_{\nu_z}p(r,x,z),
 \qquad r>0,\quad x\in\Om,\quad z\in\partial\Om.
\end{equation}
The kernel construction in \cite[Theorems~3.1, 3.12 and
Definitions~2.1, 3.3]{G} applies to the compact smooth manifold
$\overline\Om$ with its Euclidean metric.
It supplies the heat equation, initial delta and zero Dirichlet value,
smoothness up to both spatial boundaries for positive time, and
flatness of every derivative at zero on separated spatial sets.
In an interior chart it also supplies the form
$r^{-d/2}A(\sqrt r,(x-y)/\sqrt r,y)$, where $A$ is smooth and rapidly
decreasing, with every derivative, in its middle variable.
We derive the remaining assertions here.

Writing $S(r)h(x)=\int_\Om p(r,x,y)h(y)\,dy$, this local form and the initial
delta imply $S(r)h\to h$ in $L^2$ for
$h\in C^\infty_c(\Om)$.
Indeed the initial delta forces $\int A(0,X,y)\,dX=1$;
the substitution $y=x-\sqrt r X$ and rapid decay give convergence with all
derivatives on compact interior sets. Off-diagonal flatness gives convergence
to zero on a boundary collar separated from $\operatorname{supp}h$.
A finite covering gives uniform convergence on $\overline\Om$.
The energy identity and its negative-part version then give
$p\geq0$ and $P\geq0$.
The kernel is real by uniqueness of its real initial-value problem.
Differentiating a pairing of forward and backward smooth evolutions gives
symmetry; energy uniqueness gives the semigroup property.
These statements extend by density to $L^2(\Om)$.

Extend a zero-boundary smooth function to a containing cube of side
$\ell$. Integration along one coordinate and Cauchy--Schwarz give
$\|h\|_2^2\leq\ell^2\|\nabla h\|_2^2$.
The heat energy identity consequently gives
$\|S(r)\|_{2\to2}\leq e^{-r/\ell^2}$.
Splitting $S(r)=S(1)S(r-2)S(1)$ for $r>2$ and using
Cauchy--Schwarz against derivatives of the smooth time-one kernel gives
\begin{equation}\label{eq:larger}
 \sup_{x,y\in\overline\Om}
 |\partial_x^\alpha\partial_y^\beta p(r,x,y)|
          \leq C_{\alpha,\beta}e^{-(r-2)/\ell^2}.
\end{equation}
Time derivatives follow by the heat equation.

\begin{lemma}[Mass and concentration]\label{lem:mass}
For $x\in\Om$,
\begin{equation}\label{eq:mass}
 \int_0^\infty\int_{\partial\Om}P(r,x,z)\,dS(z)\,dr=1.
\end{equation}
For every fixed $\delta>0$ and all sufficiently small $\varepsilon>0$,
\begin{equation}\label{eq:concentration}
 \sup_{p\in\partial\Om}
 \int_{\{r\geq\delta\ {\rm or}\ |z-p|\geq\delta\}}
 P(r,\Psi_\varepsilon(p),z)\,dS(z)\,dr
       \leq C_\delta\varepsilon .
\end{equation}
The latter integral is over $r>0,z\in\partial\Om$.
Its constant depends only on $\Om,\delta$.
\end{lemma}
\begin{proof}
Let $M(r,x)=\int_\Om p(r,x,y)\,dy$.
An interior cutoff equal to one near $x$, the initial property and
off-diagonal flatness give $M(0+,x)=1$;
\eqref{eq:larger} gives $M(\infty,x)=0$.
Symmetry and the divergence theorem give
$\partial_rM=-\int_{\partial\Om}P(r,x,z)\,dS(z)$.
Integration and positivity prove \eqref{eq:mass}.

For $p,z\in\partial\Om$, $P(r,p,z)=0$: the kernel at the first boundary
point is identically zero as a function of its source.
The mean value formula along $p-a\nu(p)$ bounds the part $r\geq\delta$
of \eqref{eq:concentration} by $\varepsilon$ times the integral of
$\sup_{0\leq a\leq\varepsilon}|\nabla_xP(r,p-a\nu(p),z)|$.
Positive-time smoothness and \eqref{eq:larger} make this integral uniformly
finite. For $|z-p|\geq\delta$ and $\varepsilon<\delta/2$, the spatial
separation is at least $\delta/2$.
Off-diagonal flatness at zero and exponential decay at infinity make the
same derivative integral finite over all $r>0$.
Adding the two estimates proves \eqref{eq:concentration}.
\end{proof}

\subsection{Uniform \texorpdfstring{$L^2$}{L2} bounds and strong parallel traces}
Define the Poisson integral directly by
\begin{equation}\label{eq:Poisson}
 Y_g(t,x)=\int_0^t\int_{\partial\Om}P(t-s,x,z)g(s,z)\,dS(z)\,ds.
\end{equation}
To prove a uniform bound for its parallel values we need a column estimate.
Let $G(x,y)=\int_0^\infty p(r,x,y)\,dr$.
This integral is locally integrable, since its volume row integral at small
time is bounded by one and at large time decays exponentially.
Integration of the heat equation gives
$-\Delta_yG(x,y)=\delta_x(y)$, with zero boundary values.
All boundary derivatives off the interior pole can be integrated using
separation at zero and \eqref{eq:larger}. In particular
\[
 P^{\rm ell}(x,z):=-\partial_{\nu_z}G(x,z)
             =\int_0^\infty P(r,x,z)\,dr.
\]

At the uniform exterior tangent ball with center $c_z=z+R\nu(z)$,
put $a=|x-c_z|>R$ and $x^*=c_z+R^2(x-c_z)/a^2$.
Writing $\omega_{d-1}=|\mathbb S^{d-1}|$, the comparison Green function is
\begin{equation}\label{eq:ballGreen}
 H(x,y)=
 \begin{cases}
 \displaystyle\frac{|x-y|^{2-d}-(R/a)^{d-2}|x^*-y|^{2-d}}
                         {(d-2)\omega_{d-1}},&d>2,\\[5pt]
 \displaystyle\frac1{2\pi}\log\frac{a|y-x^*|}{R|y-x|},&d=2 .
 \end{cases}
\end{equation}
The identity
$a^2|y-x^*|^2-R^2|y-x|^2
=(a^2-R^2)(|y-c_z|^2-R^2)$ shows $H\geq0$ outside the ball,
and $H=0$ on its sphere.
Differentiation away from the pole and the unit radial flux of the
fundamental solution give $-\Delta_yH(x,y)=\delta_x(y)$, the same
distributional Laplacian as $G$.
Thus $H-G$ is harmonic in the bounded domain $\Om$, with nonnegative
boundary values.
For completeness, distributional harmonic regularity follows by local
radial mollification and the mean-value identity: each smooth mollification
equals its convolution with a fixed smaller radial mass-one test function;
passing to the distributional limit gives a smooth representative.
Testing the negative part proves $H-G\geq0$.
At the tangent point their common zero value and their inward derivatives,
followed by differentiation of \eqref{eq:ballGreen}, give
\begin{equation}\label{eq:ellbound}
 0\leq P^{\rm ell}(x,z)\leq
          \frac{|x-c_z|^2-R^2}{\omega_{d-1}R|x-z|^d}.
\end{equation}
The logarithmic case gives the same formula; no comparison at infinity
is needed.

By \eqref{eq:quadratic},
$|\Psi_\varepsilon(p)-z|^2\geq\frac12|p-z|^2+\varepsilon^2$
and
$|\Psi_\varepsilon(p)-c_z|^2-R^2\leq C(\varepsilon+|p-z|^2)$.
Consequently
\[
 P^{\rm ell}(\Psi_\varepsilon(p),z)
       \leq C\frac{\varepsilon+|p-z|^2}
                          {(\varepsilon^2+|p-z|^2)^{d/2}}.
\]
Finite boundary charts bound its column integral by a constant times
\[
 \int_0^{C_0}
   \frac{(\varepsilon+s^2)s^{d-2}}
                         {(\varepsilon^2+s^2)^{d/2}}\,ds .
\]
The $\varepsilon$ term is uniformly bounded after $s=\varepsilon v$.
The other term is $O(\varepsilon)$ below $\varepsilon$ and bounded above it.
The far parts of the charts are uniformly bounded.
Also \eqref{eq:ellbound} is at most $C|x-z|^{1-d}$ on bounded $\Om$,
whose volume integral is uniformly finite.
We have proved
\begin{equation}\label{eq:columns}
 \sup_{\varepsilon,z}\int_{\partial\Om}
      P^{\rm ell}(\Psi_\varepsilon(p),z)\,dS(p)<\infty,\qquad
 \sup_z\int_\Om P^{\rm ell}(x,z)\,dx<\infty .
\end{equation}

Let $T_\varepsilon g(t,p)=Y_g(t,\Psi_\varepsilon(p))$.
The nonnegative time-space kernel has row integral at most one by
\eqref{eq:mass} and uniformly bounded column integral by
\eqref{eq:columns}.
Cauchy--Schwarz with the kernel, followed by nonnegative Fubini, gives
\begin{equation}\label{eq:Schur}
 \|T_\varepsilon g\|_{L^2((0,T)\times\partial\Om)}
       \leq C^{1/2}\|g\|_2,\qquad
 \|Y_g\|_{L^2((0,T)\times\Om)}\leq C^{1/2}\|g\|_2 .
\end{equation}
Constants in these trace estimates may depend on the fixed domain and $T$
but not on the data or small parallel distance.

For $g\in C^\infty_c((0,T)\times\partial\Om)$, extend $g$ by zero in time.
Formula \eqref{eq:Poisson} is then an integral over $r>0$ of
$P(r,\Psi_\varepsilon(p),z)g(t-r,z)$.
If $\omega_g$ is a uniform modulus of continuity, the mass and
concentration estimates yield
\[
 |T_\varepsilon g(t,p)-g(t,p)|
    \leq\omega_g(2\delta)+2\|g\|_\infty C_\delta\varepsilon .
\]
First taking $\varepsilon\downarrow0$, then $\delta\downarrow0$, gives
uniform convergence.
Smooth compact-time boundary functions are dense in full boundary $L^2$:
use finitely many smooth charts and partitions, and ordinary mollification
in the chart and time variables.
Thus \eqref{eq:Schur} gives
\begin{equation}\label{eq:strongtrace}
 T_\varepsilon g\longrightarrow g
            \quad\text{in }L^2((0,T)\times\partial\Om)
 \quad\text{for every $g\in L^2$.}
\end{equation}
For example approximate by $g_k$ and bound the error by
$(C^{1/2}+1)\|g-g_k\|_2+\|T_\varepsilon g_k-g_k\|_2$.
No smoothness of a patch edge enters this density argument.

\subsection{Identification with transposition}
\begin{lemma}[Parallel trace criterion]\label{lem:trace}
Suppose $Y\in L^2((0,T)\times\Om)$ is caloric, smooth on
$[0,T]$ times each compact interior subset, and zero at $t=0$.
If $Y(t,\Psi_{\varepsilon_n}(p))\to g(t,p)$ strongly in full boundary
$L^2$ for one sequence $\varepsilon_n\downarrow0$, then $Y$ is the
zero-initial transposition state for $g$.
Conversely the state for every $g\in L^2$ is given by
\eqref{eq:Poisson}, belongs to this class and satisfies
\eqref{eq:strongtrace}.
If $g$ is supported on $\Gamma$, the state and all its derivatives are
smooth and flat at zero locally up to $\Sigma$.
\end{lemma}
\begin{proof}
Take a test $\phi$ in \eqref{eq:transposition}.
In the collar $x=p-r\nu(p)$ choose a fixed smooth cutoff $\eta(r)$
equal to one near zero and zero before leaving the collar, and on
$\Om_\varepsilon$ put
\[
 \phi_\varepsilon(t,x)=
        \phi(t,x)-\eta(r)\phi(t,\Psi_\varepsilon(p)).
\]
Outside the collar leave $\phi$ unchanged.
This test is zero on $\partial\Om_\varepsilon$ and at $T$.
Since $\phi(t,p)=0$, its correction is $O(\varepsilon)$ in
$C_x^2C_t^1$, uniformly on the closed cylinder, and
$\partial_{\nu_\varepsilon}\phi_\varepsilon(t,\Psi_\varepsilon(p))
\to\dn\phi(t,p)$ uniformly.
Integration by parts on the smooth inner domain gives
\[
 \int_0^T\int_{\Om_\varepsilon}
       Y(-\dt\phi_\varepsilon-\Delta\phi_\varepsilon)
  =-\int_0^T\int_{\partial\Om}
       Y(t,\Psi_\varepsilon(p))
       \partial_{\nu_\varepsilon}\phi_\varepsilon(t,\Psi_\varepsilon(p))
       J_\varepsilon(p)\,dS\,dt .
\]
The term involving $\partial_{\nu_\varepsilon}Y$ vanishes because the test
has zero boundary value; the initial term vanishes because $Y(0)=0$.
Strong boundary convergence and $J_\varepsilon\to1$ give the right-hand
limit along the prescribed sequence.
The disappearing collar has vanishing integral because $Y\in L^2$,
and the correction tends uniformly to zero.
This proves \eqref{eq:transposition}.

Uniqueness in that identity can be verified directly.
For $\psi\in C^\infty_c((0,T)\times\Om)$ use
\[
 \phi(t)=\int_0^{T-t}S(r)\psi(t+r)\,dr.
\]
The kernel's small-time interior convergence, separation near the boundary,
and positive-time smoothness make this test smooth on the closed cylinder.
Integrating the $r$ derivative gives $-\dt\phi-\Delta\phi=\psi$;
its boundary and terminal values vanish.
The difference of two identities therefore tests to zero against every
such $\psi$.

The directly defined $Y_g$ satisfies the volume bound and strong trace
already proved. On compact interior sets the source boundary is separated:
all differentiated kernels, extended by zero to negative time differences,
are smooth and square-integrable in the source variables.
Differentiating against $g\in L^2$ gives causal smoothness, flatness and
the heat equation. The first part of the lemma now identifies $Y_g$ with
the transposition state.
For $g=\mathbf1_\Gamma u$, the same separation argument works in a collar
of each compact subset of $\Sigma$, including the endpoints in time.
It also gives smooth convergence of all parallel normal derivatives there.
\end{proof}

\section{Causal spatial continuation}\label{sec:continuation}

All causal functions and distributions in this section are extended by zero
for negative time. Smoothness at the terminal time always means left
continuity of every derivative on compact spatial subsets.

\begin{lemma}[Local smoothing and fixed Taylor radii]\label{lem:Taylor}
Let $D$ be a bounded open set and let $F$ be a causal distributional
solution of $(\dt-\Delta)F=0$ on $(-\infty,T)\times D$ such that
$\int_0^T\int_U|F|^2\,dx\,dt<\infty$ for every $U\Subset D$.
Then $F$ is smooth on $[0,T]$ times every compact interior subset and
all its derivatives are zero at $t=0$.
For $q\in D$ put $r=\dist(q,\partial D)/(128\sqrt d)$.
Its Taylor coefficients satisfy, for every $m\in\N$,
\begin{equation}\label{eq:localbudget}
 \sum_\alpha r^{|\alpha|}
 \sup_{0\leq t\leq T}
       \left|\frac{\partial_x^\alpha\partial_t^mF(q,t)}{\alpha!}\right|
       <\infty .
\end{equation}
The series equals $F$ on the real polydisc $|x_i-q_i|<4r$.
Translation by $|\delta|<r/4$ gives its actual coefficients at $q+\delta$
by \eqref{eq:translation}, with all time derivatives converging uniformly.
\end{lemma}
\begin{proof}
For a compact spatial cutoff $\eta$ equal to one on a smaller ball,
the distributional equation gives
\[
 (\dt-\Delta)(\eta F)
       =-2\operatorname{div}(F\nabla\eta)+(\Delta\eta)F.
\]
The causal Gaussian fundamental solution consequently gives, on that ball,
\begin{equation}\label{eq:cutoff}
 F(x,t)=\int_0^t\int
 \big[-2\nabla_x\Phi(t-s,x-y)\cdot\nabla\eta(y)
              +\Phi(t-s,x-y)\Delta\eta(y)\big]F(y,s)\,dy\,ds.
\end{equation}
Here $(\dt-\Delta)\Phi=\delta_{(0,0)}$ follows by direct differentiation,
Gaussian mass one, and its approximate-identity limit.
To check uniqueness of this causal convolution identity, take the spatial
Fourier transform of the homogeneous difference, which is tempered on
each compact time interval. On a bounded frequency set multiplication
by $e^{t|\xi|^2}$ changes its equation to zero temporal derivative.
A distribution with zero temporal derivative is constant in time:
every compact time test of integral zero has a compactly supported
primitive. Zero past makes that constant zero.
Exhausting frequency space proves uniqueness. This argument uses an
exponential multiplier only on bounded frequency sets.

The derivatives of $\eta$ are separated from $x$.
Every derivative of the kernel in \eqref{eq:cutoff}, extended by zero at
nonpositive time differences, is bounded, smooth and flat at zero there.
Cauchy--Schwarz against local space-time $L^2$ justifies differentiation,
continuity up to $T$, and flatness at zero.

Now put $R=\dist(q,\partial D)/2$, choose $\eta$ supported in $B(q,R)$
and equal to one on $B(q,R/2)$, and for each $m$ set
\[
 Z_m=\eta\partial_t^mF,\qquad
 H_m=-2\nabla\eta\cdot\nabla\partial_t^mF
                   -(\Delta\eta)\partial_t^mF.
\]
The smooth bounded source $H_m$ is supported in
$R/2\leq|y-q|\leq R$ and
$Z_m=\Phi*H_m$ by the same causal uniqueness argument.
For $h=x-q$ the algebraic exponential series of the Gaussian has coefficients
$A_\alpha$, with
\[
 \sum_\alpha|A_\alpha|\rho^{|\alpha|}
 \leq(4\pi\tau)^{-d/2}
 \exp\left(-\frac{|q-y|^2}{4\tau}
       +\frac{\rho\sum_i|q_i-y_i|}{2\tau}
       +\frac{d\rho^2}{4\tau}\right).
\]
Choose $\rho=4r=R/(16\sqrt d)$.
On the source annulus the exponent is at most
$-31R^2/(1024\tau)$.
Absolute integration of this bound over $0<\tau<T$ gives
\[
 \sup_{0\leq t\leq T}
 \sum_\alpha(4r)^{|\alpha|}
  \left|\frac{\partial_x^\alpha\partial_t^mF(q,t)}{\alpha!}\right|
 \leq M_m,\qquad
 M_m=\frac{\|H_m\|_\infty|B(q,R)|}{(4\pi)^{d/2}}
     \int_0^T\tau^{-d/2}e^{-31R^2/(1024\tau)}\,d\tau .
\]
The integral is finite. This also proves that the series is the actual
function on the polydisc, which lies where $\eta=1$.
Bounding each coefficient by $M_m(4r)^{-|\alpha|}$ and summing at radius
$r$ proves \eqref{eq:localbudget}, with bound $M_m(4/3)^d$.
Differentiating the actual series at $\delta$ proves the binomial formula.
For fixed $\beta$ its binomial factors have polynomial growth in $\alpha$,
while $|\delta_i|/(4r)<1/16$; the resulting series and every time
derivative converge uniformly.
\end{proof}

\begin{lemma}[The continuation determined by the tests]\label{lem:gluing}
If tests \ref{test:A} and \ref{test:B} pass, the function $\calF_f$
in \eqref{eq:FY} is smooth and causal on $D$, caloric, and spatially
analytic at each fixed time. It equals $V_f$ on $E$.
Conversely, if $F$ satisfies the causal equation and full-interval local
$L^2$ hypothesis of Lemma~\ref{lem:Taylor} and equals $V_f$ on $E$,
the recursively computed coefficients satisfy test \ref{test:B}
and reconstruct this $F$.
\end{lemma}
\begin{proof}
On each strictly smaller covering ball a spatial derivative of the series
introduces only a polynomial in $\alpha$, absorbed by the strict reduction
in radius in \eqref{eq:budget}.
Thus every space-time derivative converges uniformly.
Flatness and the recurrence \eqref{eq:recurrence} give causality and the
heat equation term by term.
Equality on the dense rational space-time points gives equality on
every overlap, so the series glue.

The source of $V_f$ is outside $E$.
It is smooth, causal and caloric there; spatial analyticity follows either
from the separated Gaussian series or Lemma~\ref{lem:Taylor}.
The root coefficients are its actual Taylor coefficients, so the
root series agrees with $V_f$ on a neighborhood of $q_0$.
At each fixed time the analytic identity principle on connected $E$
gives equality throughout $E$.
One can see this principle directly: the set of points with a zero
neighborhood for the difference is open; at a limit point all derivatives
vanish by continuity and the local convergent series gives a zero
neighborhood, so that set is also closed.

For the converse, Lemma~\ref{lem:Taylor} applies at every center.
Along each finite tree path its translation formula recovers the actual
coefficients. The prescribed summability, flatness, heat recurrence and
overlap agreement therefore all hold.
\end{proof}

\begin{lemma}[Cancellation of the measured source]\label{lem:cancel}
For an original input $u\in L^2((0,T)\times\Gamma)$, let
$f=\dn y_u|_\Sigma$. Then
\[
 F=V_f-\mathbf1_\Om y_u
\]
is causal and caloric on $D$, equals $V_f$ on $E$, and satisfies
$\int_0^T\int_U|F|^2\,dx\,dt<\infty$ for every $U\Subset D$.
\end{lemma}
\begin{proof}
Lemma~\ref{lem:trace} gives smoothness and flatness of $y_u$ locally up to
the seed patch, where its Dirichlet value is zero.
The surface Gaussian bound
$\int_{\operatorname{supp}\chi}\Phi(r,x-z)\,dS(z)\leq Cr^{-1/2}$
follows from finite graph charts: surface measure in a radius-$a$ ball
is at most $Ca^{d-1}$; sum the Gaussian over dyadic rings of scale
$\sqrt r$. The bound is integrable in time, so $V_f$ is locally bounded
across its source surface and has zero initial value.
The fundamental-solution identity gives
$(\dt-\Delta)V_f=\chi f\,\delta_{\partial\Om}$.
Inside $D$ the only original boundary lies in $B$, where $\chi=1$.
Two integrations by parts there give
$(\dt-\Delta)(\mathbf1_\Om y_u)=f\,\delta_{\partial\Om}$,
since the zero Dirichlet value removes the derivative-of-surface term.
Explicitly, for a compact local space-time test $\varphi$,
\[
 \int_0^T\int_\Om y_u(-\dt\varphi-\Delta\varphi)\,dx\,dt
   =\int_0^T\int_{\partial\Om}
           (\dn y_u\,\varphi-y_u\,\dn\varphi)\,dS\,dt.
\]
The source sign is positive.
No initial delta occurs. Interior smoothness and flatness in the seed
collar make the zero-extended state continuous at $t=0$ in local $L^2$,
with value zero; the potential is also continuous with zero initial value.
The two surface sources therefore cancel.
Away from that boundary the assertion is immediate, and on $E$ the
extended state is zero.
\end{proof}

\section{Proof of the range and norm statements}\label{sec:rangeproof}

\begin{proof}[Proof of Theorem~\ref{thm:range}]
For necessity take any original $u\in L^2((0,T)\times\Gamma)$.
The separated representation in Lemma~\ref{lem:trace} gives test
\ref{test:A}.
Lemma~\ref{lem:cancel} supplies the causal caloric $F$ on $D$ with the
required exterior germ.
Lemma~\ref{lem:gluing} proves test \ref{test:B} and identifies
$\calY_f=y_u$.
The volume bound and strong parallel trace in Lemma~\ref{lem:trace}
give test \ref{test:C} and
$G_n(f)\to\mathbf1_\Gamma u$ in full boundary $L^2$.
This gives test \ref{test:D}.
Separation of the input and observation closures gives smooth convergence
of the parallel normal derivatives on every compact subset of $\Sigma$,
which is test \ref{test:E}.

For sufficiency let all five tests pass.
Lemma~\ref{lem:gluing} makes $\calF_f$ smooth, causal and caloric on $D$.
Since the source of $V_f$ lies on the boundary, $\calY_f$ is smooth and
caloric inside $\Om$, with zero initial value.
By test \ref{test:C}, it is in volume $L^2$ and its parallel traces
converge strongly to some $g$ in full boundary $L^2$.
Lemma~\ref{lem:trace} gives the original transposition identity for $\calY_f$.
Test \ref{test:D} makes $g=0$ almost everywhere on
$\partial\Om\setminus\Gamma$, so $u=g|_\Gamma$ is an admissible input.
Uniqueness identifies $\calY_f=y_u$.
This original state is smooth locally up to $\Sigma$ by separation.
Its actual normal derivative is therefore the parallel limit in test
\ref{test:E}, namely $f$ on the entire original observation patch.

If two inputs give the same flux, their continued functions from
Lemma~\ref{lem:cancel} agree on $E$.
They are spatially analytic at each time by Lemma~\ref{lem:Taylor}.
The identity principle on connected $D$ gives equality everywhere,
then equality of states on $\Om$ and of their parallel $L^2$ traces.
Thus the two inputs are equal.
Strong trace convergence proves \eqref{eq:N}; absolute continuity of the
integral makes $N_f(t)$ finite and continuous, with $N_f(0)=0$.
The range map is linear, bijective and isometric in $N_f(T)$.
Transporting the $L^2$ inner product gives completeness and the Hilbert
structure.
Finally the original range and the unique input do not refer to any
auxiliary choices, proving invariance of the class and norm.
\end{proof}

\section{Sharp restrictions at the initial time}\label{sec:sharp}

For an observation point $x\in\Sigma$ define the forward flux kernel and
its evaluation scale by
\begin{equation}\label{eq:KC}
 K(r,x,z)=-\partial_{\nu_x}\partial_{\nu_z}p(r,x,z),\qquad
 C_{\Gamma,x}(t)^2=\int_0^t\int_\Gamma|K(r,x,z)|^2\,dS(z)\,dr .
\end{equation}
These are computed from the Dirichlet heat kernel of the fixed domain,
not from a tracking inverse. The derivatives are taken from the interior.
Separation makes $C_{\Gamma,x}(t)$ finite; its square is smooth and
flat at zero. Consequently $C_{\Gamma,x}(t)=O(t^M)$ for every $M>0$.
We do not need any assertion about derivatives of the square root at zero.

\begin{theorem}[Exact evaluation scale]\label{thm:sharp}
Under the hypotheses of Theorem~\ref{thm:range}, fix $x\in\Sigma$.
For every $0<t\leq T$, $C_{\Gamma,x}(t)>0$ and
\begin{equation}\label{eq:sharpbound}
 |f(t,x)|\leq C_{\Gamma,x}(t)N_f(t)
                  \quad(f\in\calX_T),\qquad
 \frac{|f(t,x)|}{C_{\Gamma,x}(t)}\longrightarrow0
                         \quad(t\downarrow0).
\end{equation}
The scale $C_{\Gamma,x}(t)$ is exactly the norm of evaluation at $(t,x)$
on the unit ball for $N_f(T)$; the constant one is optimal.
For every function $\eta(t)>0$ with $\eta(t)\to0$ as $t\downarrow0$
there is one fixed $f\in\calX_T$ such that
\begin{equation}\label{eq:norate}
 \sup_{0<t<T}\frac{|f(t,x)|}{C_{\Gamma,x}(t)\eta(t)}=\infty .
\end{equation}
Thus the fixed-target little-o assertion has no prescribed extra rate
valid for all targets.
\end{theorem}

This gives an exact quantitative restriction including the full
geometry-dependent scale at small times.
The pointwise bound alone is not asserted to be sufficient for tracking;
the continuation, trace and patch compatibility tests in
Section~\ref{sec:class} give the necessary and sufficient class.

\begin{lemma}[Strict kernel signs]\label{lem:sign}
For $r>0$, $w,y\in\Om$ and distinct boundary points $x,z$,
\[
 p(r,w,y)>0,\qquad P(r,w,z)>0,\qquad K(r,x,z)<0 .
\]
\end{lemma}
\begin{proof}
We give elementary positivity propagation and a boundary barrier.
In the cube $(-a,a)^d$ let
$\varphi(X)=\prod_i\cos(\pi X_i/(2a))$ and
$\lambda=d\pi^2/(4a^2)$.
If its center moves as $c(s)=c_0+v(s-s_0)$, the function
\[
 b(s,w)=b_0e^{-v\cdot(w-c(s))/2}
 e^{-(\lambda+|v|^2/4)(s-s_0)}\varphi(w-c(s))
\]
is caloric, positive inside and zero on the moving boundary.
If a nonnegative caloric $h$ is positive near the initial closed cube,
choose $b_0>0$ small enough that $b(s_0,\cdot)\leq h(s_0,\cdot)$.
When the moving cube stays in $\Om$, comparison gives $h\geq b$:
in moving coordinates the squared negative-part energy is nonincreasing,
and the constant transport term integrates to zero because the negative
part has zero boundary trace.
A polygonal path compactly contained in $\Om$ can be traversed by small
moving cubes in any prescribed positive time.
At each junction shrink the next initial cube inside the previous
terminal open cube. Thus positivity on one open set propagates to every
interior point in any additional positive time.
For fixed source $y$, the nonnegative kernel's initial delta implies a
positive open set at an arbitrarily small positive time.
Propagation proves strict interior positivity.

At a boundary point take an interior tangent ball $B(c,R)$ and an interval
$[r-\delta,r]$, $0<\delta<r$.
In its shell $R/2<|w-c|<R$ put
\[
 b(w,s)=e^{-A(|w-c|^2+\mu(r-s))}-e^{-AR^2},
                 \qquad \mu\delta\geq3R^2/4.
\]
This function is nonpositive initially and on the outer sphere.
A positive caloric function has a positive minimum on the compact inner
sphere over this time interval, so a small positive multiple of $b$
lies below it there.
For sufficiently large $A$,
\[
 (\partial_s-\Delta)b
   =A(\mu+2d-4A|w-c|^2)e^{-A(|w-c|^2+\mu(r-s))}\leq0.
\]
Negative-part comparison in the fixed shell gives the lower bound.
At the final tangency point the inward derivative is
$2AR e^{-AR^2}>0$, so the outward derivative of a positive caloric
function with zero boundary value is strictly negative.
Apply this first to the kernel in its source variable, using symmetry,
to obtain $P>0$.
For fixed $z$, $P$ is positive and caloric in the interior variable,
and zero on that boundary at every positive time.
A second application gives $\partial_{\nu_x}P=K<0$.
\end{proof}

\begin{proof}[Proof of Theorem~\ref{thm:sharp}]
Since $\Gamma$ is nonempty and relatively open, it has positive surface
measure; Lemma~\ref{lem:sign} proves $C_{\Gamma,x}(t)>0$.
The separated Poisson representation and the recovered input give
\[
 f(t,x)=\int_0^t\int_\Gamma K(t-s,x,z)u(s,z)\,dS(z)\,ds .
\]
Cauchy--Schwarz and \eqref{eq:N} prove \eqref{eq:sharpbound}.
For the second assertion $N_f(t)\to0$ by absolute continuity of the
squared-input integral.
At a fixed positive $t$, the allowed input
$u_t(s,z)=\overline{K(t-s,x,z)}/C_{\Gamma,x}(t)$ on
$(0,t)\times\Gamma$, zero afterwards, has total norm one and output
$f_t(t,x)=C_{\Gamma,x}(t)$.
The norm equality in Theorem~\ref{thm:range} proves exact optimality.

It remains to show \eqref{eq:norate} for one fixed input.
Write $C(t)=C_{\Gamma,x}(t)$ and choose $t_n\downarrow0$ recursively with
$t_{n+1}<t_n/2$, $\eta(t_n)\leq n^{-2}$, and
\[
 \|\mathbf1_{(0,t_{n+1})}(s)K(t_n-s,x,z)\|_{L^2(ds\,dS)}
                  \leq C(t_n)/(8n).
\]
After $t_n$ is fixed, absolute continuity of its squared kernel integral
permits the next upper bound on $t_{n+1}$, and the limit of $\eta$ permits
its next smallness condition.
On disjoint intervals $I_n=(t_{n+1},t_n)$ let $e_n$ be the conjugated
kernel $K(t_n-s,x,z)$ restricted to $I_n\times\Gamma$ and divided by its
$L^2$ norm.
This norm is at least $C(t_n)\sqrt{1-(8n)^{-2}}$, so $e_n$ is well defined
with unit norm.
The fixed input $u=\sum_{n\geq1}n^{-1}e_n$ has
$\|u\|_2^2=\sum_{n\geq1}n^{-2}<4$.
At $t_n$ earlier indices are in the future and contribute nothing.
The self term contributes at least
$C(t_n)n^{-1}\sqrt{1-(8n)^{-2}}$.
All later terms have norm below two, are supported in $(0,t_{n+1})$,
and contribute in absolute value at most $C(t_n)/(4n)$.
Hence $|f(t_n,x)|\geq C(t_n)/(2n)$ and the quotient in
\eqref{eq:norate} is at least $n/2$.
Theorem~\ref{thm:range} makes this one output a member of the complete
target class, proving the assertion.
\end{proof}

\paragraph{AI-use statement.}
OpenAI Codex assisted with literature retrieval, mathematical exploration,
proof development, checking, and preparation of this manuscript.
The submission is a mathematical claim for external review.

\begin{thebibliography}{9}
\bibitem{BZ}
J. A. B\'arcena-Petisco and E. Zuazua,
\emph{Tracking controllability for the heat equation},
IEEE Transactions on Automatic Control \textbf{70}(3) (2025), 1935--1940.
\href{https://doi.org/10.1109/TAC.2024.3476174}{doi:10.1109/TAC.2024.3476174}.
Author manuscript:
\href{https://arxiv.org/abs/2310.00314v3}{arXiv:2310.00314v3}.

\bibitem{DL}
L. Davron and P. Lissy,
\emph{Exact output tracking for the one-dimensional heat equation and
applications to the interpolation problem in Gevrey classes of order 2},
preprint, 2025.
\href{https://arxiv.org/abs/2511.07326v2}{arXiv:2511.07326v2}.

\bibitem{PS}
R. E. Puzyrev and A. A. Shlapunov,
\emph{On an Ill-Posed Problem for the Heat Equation},
Journal of Siberian Federal University. Mathematics \& Physics
\textbf{5}(3) (2012), 337--348.
\href{https://elib.sfu-kras.ru/bitstream/2311/2924/1/shlapunev.pdf}{Authoritative full text}.

\bibitem{G}
D. Grieser,
\emph{Notes on Heat Kernel Asymptotics},
13 August 2004.
\href{https://web.math.ku.dk/grubb/notes/heat.pdf}{Full text}.

\bibitem{HM}
J. K. Hunter,
\emph{Measure Theory}, 2011.
\href{https://www.math.ucdavis.edu/~hunter/measure_theory/measure_notes.pdf}{Lecture notes}.

\bibitem{HP}
J. K. Hunter,
\emph{Notes on Partial Differential Equations}, revised 18 June 2014.
\href{https://www.math.ucdavis.edu/~hunter/pdes/pde_notes.pdf}{Lecture notes}.
\end{thebibliography}
\end{document}
